\begin{corollary}[Inertia group constraint of central six-fold aggregation: Perron branch algebraic degree must be divisible by \(6\)]\label{cor:rate-center-perron-degree-multiple-6}
Adopting the notation of the elimination certificate \(R(\tfrac12,u)=-(u-1)^6Q(u)/64\), and assuming the generic transversality condition
\[
\partial_\alpha R\!\left(\tfrac12,1\right)\neq 0.
\]
Let \(F(\alpha,u)\in\QQ[\alpha,u]\) be the irreducible factor of \(R(\alpha,u)\) in \(\QQ(\alpha)[u]\) that contains the point \((\alpha,u)=(\tfrac12,1)\) (for instance, the factor containing the Perron branch).
Then the ramification index at \(\alpha=\tfrac12\) of the function field extension
\(
\QQ(\alpha)\subset \QQ(\alpha,u)\ (F(\alpha,u)=0)
\)
is \(e=6\);
thus the \(u\)-degree of \(F\) (i.e., the algebraic degree of this branch over \(\QQ(\alpha)\)) must be divisible by \(6\).
\end{corollary}
\begin{proof}
By the certificate \(R(\tfrac12,u)=-(u-1)^6Q(u)/64\) with \(Q(1)\neq 0\), on the central section \(\alpha=\tfrac12\), the point \(u=1\) is a root of exact multiplicity \(6\).
Under the transversality condition \(\partial_\alpha R(\tfrac12,1)\neq 0\), the Weierstrass preparation theorem (or equivalent formulation of algebraic ramification theory) gives: in some neighborhood of \((\tfrac12,1)\), the local solutions of the equation \(R(\alpha,u)=0\) satisfy
\(
u-1\asymp (\alpha-\tfrac12)^{1/6}
\),
i.e., the local ramification index at this point is \(e=6\) (see the preceding order-\(6\) singular inversion scaling).
Therefore, in the function field extension defined by the irreducible factor \(F\), the ramification index of the prime divisor corresponding to this point is \(6\);
and in algebraic function fields in characteristic \(0\), any ramification index \(e\) must divide the extension degree \([\QQ(\alpha,u):\QQ(\alpha)]\), thus \(\deg_u F\) must be divisible by \(6\). This completes the proof.
\end{proof}

\begin{remark}[Order-\(6\) cyclic inertia in Galois closure (semantic alignment)]\label{rem:rate-center-inertia-c6}
In characteristic \(0\), the local inertia group is cyclic with order equal to the ramification index.
Therefore, the central six-fold aggregation described in Corollary \ref{cor:rate-center-perron-degree-multiple-6} forces the appearance of a cyclic inertia subgroup of order \(6\) in the corresponding Galois closure;
this constraint is completely forced by the auditable factor \((u-1)^6\) of the central aggregation and does not depend on additional numerical fitting.
This remark is used only to align the "inertia/ramification" semantics and is not a prerequisite for other proofs.
\end{remark}
